\paragraph{Zeckendorf--Gödel 素数码集合的密度、Lee--Yang 实根结构与谱行列式重整}
设 \((p_k)_{k\ge 1}\) 为递增素数序列，并定义
\[
\mathrm{ZG}
:=\Bigl\{
n\in\NN:\ \forall p,\ v_p(n)\in\{0,1\},\ \forall k\ge 1,\ v_{p_k}(n)\,v_{p_{k+1}}(n)=0
\Bigr\}.
\]
即 \(n\) 为平方因子自由，且其素因子索引集合不含相邻整数。

\begin{theorem}[Gödel--Zeckendorf 素数约束 Dirichlet 级数的对数重整化与半临界解析性]\label{thm:xi-zg-dirichlet-log-renormalization}
\leanverified{paper\_xi\_zg\_dirichlet\_log\_renormalization}
令
\[
\Omega:=\Bigl\{\varepsilon=(\varepsilon_k)_{k\ge 1}:\ \varepsilon_k\in\{0,1\},\ \varepsilon_k\varepsilon_{k+1}=0,\ \varepsilon\ \text{仅有限支撑}\Bigr\}.
\]
对 \(\Re(s)>1\) 定义 Dirichlet 级数
\[
F(s):=\sum_{n\in \mathrm{ZG}}n^{-s}
=\sum_{\varepsilon\in\Omega}\ \prod_{k\ge 1}p_k^{-s\varepsilon_k}.
\]
并对 \(N\ge 0\) 定义截断
\[
F_N(s):=\sum_{\substack{\varepsilon\in\{0,1\}^N\\ \varepsilon_k\varepsilon_{k+1}=0}}\ \prod_{k=1}^{N}p_k^{-s\varepsilon_k},
\qquad
F_0(s)=1,\ \ F_1(s)=1+2^{-s}.
\]
令 \(r_N(s):=F_N(s)/F_{N-1}(s)\)（\(N\ge 1\)），并令
\[
\Delta_N(s):=\log r_N(s)-p_N^{-s}\qquad(N\ge 2),
\]
其中对数取为与实轴上 \(s>1\) 的正值一致的解析分支。
则：
\begin{enumerate}
  \item 对所有 \(\Re(s)>0\) 成立递推恒等式
  \[
  F_N(s)=F_{N-1}(s)+p_N^{-s}F_{N-2}(s),\qquad
  r_N(s)=1+\frac{p_N^{-s}}{r_{N-1}(s)}.
  \]
  \item 级数
  \[
  G(s):=\sum_{N\ge 2}\Delta_N(s)
  \]
  在每个半平面 \(\Re(s)\ge \sigma_0>\tfrac12\) 上绝对且一致收敛，从而在 \(\Re(s)>\tfrac12\) 上给出全纯函数。
  \item 令素数 \(\zeta\) 函数（prime zeta）
  \(
  P(s):=\sum_{k\ge 1}p_k^{-s}
  \)
  （在 \(\Re(s)>1\) 绝对收敛）。
  则对 \(\Re(s)>1\) 有严格分解
  \[
  \log F(s)=P(s)+G(s),
  \qquad\text{从而}\qquad
  F(s)=\exp\!\bigl(P(s)+G(s)\bigr).
  \]
\end{enumerate}
\end{theorem}

\begin{proof}
第一条由末位 \(\varepsilon_N\in\{0,1\}\) 分拆：\(\varepsilon_N=0\) 贡献为 \(F_{N-1}\)，而 \(\varepsilon_N=1\) 强制 \(\varepsilon_{N-1}=0\) 并贡献 \(p_N^{-s}F_{N-2}\)。
对 \(r_N=F_N/F_{N-1}\) 取比即得递推。
\par
对第二条，令 \(\sigma=\Re(s)\ge \sigma_0>\tfrac12\)。
由递推 \(r_N=1+p_N^{-s}/r_{N-1}\) 与 \(|p_N^{-s}|\le p_N^{-\sigma}\) 可在每个紧子域上固定 \(r_{N-1}(s)\) 远离 \(0\) 的一致下界，
从而存在常数 \(C(\sigma_0)\) 使
\[
|\Delta_N(s)|
\le
C(\sigma_0)\Bigl(p_N^{-2\sigma}+(p_Np_{N-1})^{-\sigma}\Bigr)
\qquad(\sigma\ge\sigma_0).
\]
由于
\(
\sum_N p_N^{-2\sigma}
\)
与
\(
\sum_N(p_Np_{N-1})^{-\sigma}
\)
在 \(\sigma>\tfrac12\) 收敛，故 \(\sum_{N\ge 2}\Delta_N(s)\) 在 \(\Re(s)\ge\sigma_0\) 上绝对一致收敛，从而 \(G\) 在 \(\Re(s)>\tfrac12\) 全纯。
\par
第三条由
\(
\log F_N(s)=\sum_{k=1}^N\log r_k(s)
\)
与 \(\Delta_k=\log r_k-p_k^{-s}\) 得
\[
\log F_N(s)=\sum_{k=1}^{N}p_k^{-s}+\sum_{k=2}^{N}\Delta_k(s).
\]
令 \(N\to\infty\) 并在 \(\Re(s)>1\) 用绝对收敛交换极限，得到 \(\log F(s)=P(s)+G(s)\)。证毕。
\end{proof}

\begin{theorem}[平方自由 Euler 因子分离与 \texorpdfstring{$\zeta$}{zeta}-因子化]\label{thm:xi-zg-zeta-factorization}
\leanverified{paper\_xi\_zg\_zeta\_factorization}
沿用定理 \ref{thm:xi-zg-dirichlet-log-renormalization} 的记号，并令
\[
B_N(s):=\prod_{k=1}^{N}\bigl(1+p_k^{-s}\bigr),\qquad
H_N(s):=\frac{F_N(s)}{B_N(s)}.
\]
则存在唯一函数 \(H\) 满足：
\begin{enumerate}
  \item \(H_N(s)\to H(s)\) 在每个紧集 \(K\subset\{s:\Re(s)>\tfrac12\}\) 上一致收敛；
  \item \(H(s)\) 在半平面 \(\Re(s)>\tfrac12\) 上全纯且处处非零；
  \item 在 \(\Re(s)>1\) 上成立严格因子化
  \[
  F(s)=H(s)\frac{\zeta(s)}{\zeta(2s)}.
  \]
\end{enumerate}
从而 \(F\) 在 \(\Re(s)>\tfrac12\) 上亚纯，且在该半平面内唯一可能的极点位于 \(s=1\)，并且为一阶极点。
\end{theorem}

\begin{proof}
由递推可写 \(F_N(s)=\prod_{k=1}^{N}r_k(s)\)，从而
\[
\log H_N(s)
=\sum_{k=1}^{N}\Bigl(\log r_k(s)-\log\bigl(1+p_k^{-s}\bigr)\Bigr).
\]
对 \(k\ge 2\) 使用 \(\Delta_k(s)=\log r_k(s)-p_k^{-s}\)，并将上式改写为
\[
\log H_N(s)
=\Bigl(\log r_1(s)-\log(1+p_1^{-s})\Bigr)
\sum_{k=2}^{N}\Delta_k(s)
\sum_{k=2}^{N}\Bigl(p_k^{-s}-\log(1+p_k^{-s})\Bigr).
\]
第一项恒为 \(0\)。
由定理 \ref{thm:xi-zg-dirichlet-log-renormalization}，\(\sum_{k\ge 2}\Delta_k(s)\) 在 \(\Re(s)>\tfrac12\) 的每个紧集上一致收敛。
另一方面，当 \(\Re(s)\ge \sigma_0>\tfrac12\) 时有
\(
\bigl|p_k^{-s}-\log(1+p_k^{-s})\bigr|
\ll p_k^{-2\sigma_0}
\),
从而 \(\sum_{k\ge 2}\bigl(p_k^{-s}-\log(1+p_k^{-s})\bigr)\) 在同一紧集上一致绝对收敛。
故 \(\log H_N\) 在 \(\Re(s)>\tfrac12\) 的紧集上一致收敛到某个全纯函数 \(\log H\)，并且 \(H=\exp(\log H)\) 全纯且不取零值。

对 \(\Re(s)>1\)，由平方自由 Dirichlet 级数的标准恒等式
\[
\prod_{p}\bigl(1+p^{-s}\bigr)
=\sum_{\substack{m\ge 1\\ m\ \text{平方因子自由}}}\frac1{m^s}
=\frac{\zeta(s)}{\zeta(2s)}
\]
（参见 \cite{Apostol1976,Titchmarsh1986RiemannZeta}），得 \(B_N(s)\to \zeta(s)/\zeta(2s)\)。
再由 \(F_N(s)=H_N(s)B_N(s)\) 并令 \(N\to\infty\)（在 \(\Re(s)>1\) 可由绝对收敛交换极限），得到 \(F(s)=H(s)\zeta(s)/\zeta(2s)\)。
最后，\(\zeta(2s)\) 在 \(\Re(s)>\tfrac12\) 上解析且无零点，而 \(\zeta(s)\) 在该半平面内唯一奇性为 \(s=1\) 的一阶极点，因而 \(F\) 在 \(\Re(s)>\tfrac12\) 上仅可能在 \(s=1\) 处有一阶极点。证毕。
\end{proof}

在平方自由 Euler 因子分离之外，\(\mathrm{ZG}\) 的 Dirichlet 级数还可进一步写成一条递归重整化的乘法链，其局部因子由相邻素数禁配约束逐步修正。

\begin{theorem}[Zeckendorf--Gödel Dirichlet 级数的 continued-Euler 型分解]\label{thm:xi-zg-continued-euler-factorization}
\leanverified{paper\_xi\_zg\_continued\_euler\_factorization}
沿用定理 \ref{thm:xi-zg-dirichlet-log-renormalization} 的记号。对实数 \(s>1\)，定义递归因子
\[
R_1(s):=p_1^{-s},
\qquad
R_n(s):=\frac{p_n^{-s}}{1+R_{n-1}(s)}
\qquad(n\ge 2).
\]
则对每个 \(n\ge 1\) 都有
\[
\boxed{\;
0<R_n(s)<p_n^{-s}.\;}
\]
并且 \(\mathrm{ZG}\) 的 Dirichlet 级数
\[
F(s)=\sum_{n\in \mathrm{ZG}}n^{-s}
\]
具有精确乘法分解
\[
\boxed{\;
F(s)=\prod_{n=1}^{\infty}\bigl(1+R_n(s)\bigr).\;}
\]
特别地，
\[
\boxed{\;
1<F(s)<\prod_{n=1}^{\infty}\bigl(1+p_n^{-s}\bigr)=\frac{\zeta(s)}{\zeta(2s)}
\qquad(s>1).\;}
\]
此外，\(F(s)\) 在区间 \((1,\infty)\) 上严格递减，并满足
\[
\boxed{\;
\lim_{s\to+\infty}F(s)=1.\;}
\]
\end{theorem}
\begin{proof}
由定理 \ref{thm:xi-zg-dirichlet-log-renormalization}，截断和
\[
F_N(s):=\sum_{\substack{\varepsilon\in\{0,1\}^N\\ \varepsilon_k\varepsilon_{k+1}=0}}
\prod_{k=1}^{N}p_k^{-s\varepsilon_k}
\]
满足递推
\[
F_N(s)=F_{N-1}(s)+p_N^{-s}F_{N-2}(s),
\qquad
F_0(s)=1,\quad F_1(s)=1+p_1^{-s}.
\]
再记
\[
r_N(s):=\frac{F_N(s)}{F_{N-1}(s)}
\qquad(N\ge 1).
\]
则同一定理给出
\[
r_1(s)=1+p_1^{-s},
\qquad
r_N(s)=1+\frac{p_N^{-s}}{r_{N-1}(s)}
\qquad(N\ge 2).
\]
定义
\[
R_N(s):=r_N(s)-1.
\]
于是
\[
R_1(s)=p_1^{-s},
\qquad
R_N(s)=\frac{p_N^{-s}}{1+R_{N-1}(s)}
\qquad(N\ge 2),
\]
即得所示递推。由 \(R_1(s)>0\) 出发归纳可知 \(R_N(s)>0\)；再因 \(1+R_{N-1}(s)>1\)，便有
\[
R_N(s)<p_N^{-s}.
\]

另一方面，
\[
F_N(s)=F_{N-1}(s)\,r_N(s)=F_{N-1}(s)\bigl(1+R_N(s)\bigr),
\]
故迭代得到
\[
F_N(s)=\prod_{n=1}^{N}\bigl(1+R_n(s)\bigr).
\]
由于 \(s>1\) 时 \(F_N(s)\to F(s)\)，令 \(N\to\infty\) 即得
\[
F(s)=\prod_{n=1}^{\infty}\bigl(1+R_n(s)\bigr).
\]

由 \(0<R_n(s)<p_n^{-s}\)，逐项比较可得
\[
1+R_n(s)<1+p_n^{-s}.
\]
因此
\[
1<F(s)=\prod_{n\ge 1}\bigl(1+R_n(s)\bigr)
<
\prod_{n\ge 1}\bigl(1+p_n^{-s}\bigr).
\]
最后一项等于 \(\zeta(s)/\zeta(2s)\) 是定理 \ref{thm:xi-zg-zeta-factorization} 中已经使用的平方自由 Euler 乘积恒等式。

再证单调性。若 \(1<s_1<s_2\)，则对每个 \(n\in\mathrm{ZG}\) 有
\[
n^{-s_1}\ge n^{-s_2},
\]
且在 \(n=2\) 处严格不等，于是
\[
F(s_1)-F(s_2)
=
\sum_{n\in\mathrm{ZG}}\bigl(n^{-s_1}-n^{-s_2}\bigr)>0.
\]
故 \(F\) 在 \((1,\infty)\) 上严格递减。

最后，由上界
\[
F(s)<\frac{\zeta(s)}{\zeta(2s)}
\]
且
\[
\frac{\zeta(s)}{\zeta(2s)}\longrightarrow 1
\qquad(s\to+\infty),
\]
并结合显然的下界 \(F(s)>1\)，由夹逼定理得到
\[
\lim_{s\to+\infty}F(s)=1.
\]
证毕。
\end{proof}


\begin{corollary}[素数路径硬核常数与 \texorpdfstring{$s=1$}{s=1} 留数]\label{cor:xi-zg-hardcore-constant-residue}
在定理 \ref{thm:xi-zg-zeta-factorization} 的记号下，记
\[
\mathfrak C_{\mathrm{HC}}:=H(1)=\lim_{N\to\infty}\frac{F_N(1)}{\prod_{k=1}^{N}\bigl(1+\frac1{p_k}\bigr)}\in(0,1).
\]
则 \(F(s)\) 在 \(s=1\) 的留数满足
\[
\operatorname*{Res}_{s=1}F(s)=\frac{\mathfrak C_{\mathrm{HC}}}{\zeta(2)}=\frac{6}{\pi^2}\,\mathfrak C_{\mathrm{HC}}.
\]
并且该留数与 \(\mathrm{ZG}\) 的 Dirichlet 密度一致；结合定理 \ref{con:xi-zg-density-closed-form-tail-bound} 的自然密度存在性，得到
\[
\delta_{\mathrm{ZG}}=\frac{\mathfrak C_{\mathrm{HC}}}{\zeta(2)}.
\]
数值核验可得到
\[
\mathfrak C_{\mathrm{HC}}=0.8461958103\ldots,
\qquad
\delta_{\mathrm{ZG}}=0.5144253666\ldots
\]
\leanverified{paper\_xi\_zg\_hardcore\_constant\_residue}
\end{corollary}

\begin{proof}
由定理 \ref{thm:xi-zg-zeta-factorization}，在 \(s=1\) 的邻域内
\(
F(s)=H(s)\zeta(s)/\zeta(2s)
\)
且 \(H\) 在 \(s=1\) 解析。
因此
\[
\operatorname*{Res}_{s=1}F(s)
=H(1)\operatorname*{Res}_{s=1}\frac{\zeta(s)}{\zeta(2s)}
=\frac{H(1)}{\zeta(2)}.
\]
留数等于 \(\mathrm{ZG}\) 的 Dirichlet 密度是推论 \ref{cor:xi-zg-dirichlet-residue} 的内容，而该推论亦指出留数与自然密度一致，从而得 \(\delta_{\mathrm{ZG}}=\mathfrak C_{\mathrm{HC}}/\zeta(2)\)。
最后，由于对实数 \(s>1\) 有 \(F_N(s)\le B_N(s)\) 且约束非平凡，故 \(0<H(s)<1\) 并推出 \(\mathfrak C_{\mathrm{HC}}\in(0,1)\)。证毕。
\end{proof}

\begin{proposition}[\texorpdfstring{$s=1$}{s=1} 的 Laurent 展开与 Euler--Kronecker 型常数]\label{prop:xi-zg-hardcore-laurent}
记
\[
\kappa_{\mathrm{HC}}:=\left.\frac{d}{ds}\log H(s)\right|_{s=1}.
\]
则 \(\kappa_{\mathrm{HC}}\) 存在，并且 \(F\) 在 \(s=1\) 的 Laurent 展开为
\[
F(s)=\frac{\mathfrak C_{\mathrm{HC}}}{\zeta(2)}\cdot\frac{1}{s-1}
\;+\;
\frac{\mathfrak C_{\mathrm{HC}}}{\zeta(2)}
\left(\gamma-\frac{2\zeta'(2)}{\zeta(2)}+\kappa_{\mathrm{HC}}\right)
 + O(s-1),
\qquad s\to 1,
\]
其中 \(\gamma\) 为 Euler--Mascheroni 常数。
数值核验表明 \(\kappa_{\mathrm{HC}}=0.40404\ldots\)。
\end{proposition}

\begin{proof}
由定理 \ref{thm:xi-zg-zeta-factorization} 的构造，\(\log H_N(s)\) 在 \(\Re(s)>\tfrac12\) 的紧集上一致收敛为 \(\log H(s)\)，并且其尾项由 \(\sum_{k\ge 2}\Delta_k(s)\) 与 \(\sum_{k\ge 2}(p_k^{-s}-\log(1+p_k^{-s}))\) 的一致收敛控制。
同样地，对 \(s\) 位于 \(s=1\) 的小邻域且满足 \(\Re(s)>\tfrac12\)，上述两级数可逐项求导，从而 \(\kappa_{\mathrm{HC}}\) 存在。

另一方面，
\(
\zeta(s)=\frac1{s-1}+\gamma+O(s-1)
\)
且
\[
\frac{1}{\zeta(2s)}
\;=\;
\frac{1}{\zeta(2)}
\left(1-\frac{2\zeta'(2)}{\zeta(2)}(s-1)+O((s-1)^2)\right).
\]
再由
\(
H(s)=H(1)\bigl(1+\kappa_{\mathrm{HC}}(s-1)+O((s-1)^2)\bigr)
\)
三者相乘即得所述展开式。证毕。
\end{proof}

% Auto-split to keep each TeX source < 600 lines.

\begin{theorem}[Zeckendorf--Gödel 素数码集合计数函数的幂节省误差]\label{thm:xi-zg-counting-power-saving-error}
记
\[
A(x):=\#\{n\le x:\ n\in\mathrm{ZG}\}
=\sum_{n\le x}\mathbf 1_{\mathrm{ZG}}(n).
\]
则对任意 \(\varepsilon>0\)，当 \(x\to\infty\) 有
\[
A(x)=\delta_{\mathrm{ZG}}\,x+O_\varepsilon\!\bigl(x^{\frac12+\varepsilon}\bigr).
\]
\end{theorem}

\begin{proof}
令 \(F(s)=\sum_{n\in\mathrm{ZG}}n^{-s}\)。
由定理 \ref{thm:xi-zg-zeta-factorization}，\(F\) 在 \(\Re(s)>\tfrac12\) 上亚纯且在 \(s=1\) 处有一阶极点；
其留数为 \(\delta_{\mathrm{ZG}}\)（推论 \ref{cor:xi-zg-hardcore-constant-residue}）。
\par
对非负系数 Dirichlet 级数应用 Perron 反演（参见 \cite{Apostol1976,Titchmarsh1986RiemannZeta}）可表示 \(A(x)\) 为沿竖线 \(\Re(s)=c>1\) 的积分。
将积分线移至 \(\Re(s)=\tfrac12+\varepsilon\) 时穿越 \(s=1\) 的留数贡献为 \(\delta_{\mathrm{ZG}}x\)。
由于 \(\zeta(2s)\) 在 \(\Re(s)\ge\tfrac12+\varepsilon\) 上解析且远离 \(0\)，并且 \(H(s)\) 在该半平面解析且由其对数展开的绝对收敛给出多项式增长界（定理 \ref{thm:xi-zg-dirichlet-log-renormalization}--\ref{thm:xi-zg-zeta-factorization}），结合 \(\zeta(s)\) 的经典竖线界即可控制移线后的积分余项，从而得到
\(
A(x)=\delta_{\mathrm{ZG}}x+O_\varepsilon(x^{\frac12+\varepsilon})
\)。
证毕。
\end{proof}

\begin{corollary}[Zeckendorf--Gödel 素数码集合的倒数调和渐近]\label{cor:xi-zg-reciprocal-harmonic-asymptotic}
在定理 \ref{thm:xi-zg-counting-power-saving-error} 的记号下，对任意 \(0<\varepsilon<\tfrac12\)，存在常数 \(C_{\mathrm{ZG}}\in\RR\) 使
\[
\sum_{\substack{n\le X\\ n\in\mathrm{ZG}}}\frac{1}{n}
=\delta_{\mathrm{ZG}}\log X+C_{\mathrm{ZG}}+O_\varepsilon\!\bigl(X^{-\frac12+\varepsilon}\bigr)
\qquad(X\to\infty).
\]
特别地，
\[
\sum_{\substack{n\le X\\ n\in\mathrm{ZG}}}\frac{1}{n}
\sim \delta_{\mathrm{ZG}}\log X.
\]
\end{corollary}

\begin{proof}
记误差项
\[
E(x):=A(x)-\delta_{\mathrm{ZG}}x.
\]
由定理 \ref{thm:xi-zg-counting-power-saving-error}，对任意 \(0<\varepsilon<\tfrac12\) 有
\[
E(x)=O_\varepsilon\!\bigl(x^{\frac12+\varepsilon}\bigr).
\]
对计数函数 \(A\) 应用 Abel 求和，得到
\[
\sum_{\substack{n\le X\\ n\in\mathrm{ZG}}}\frac{1}{n}
=\frac{A(X)}{X}+\int_{1}^{X}\frac{A(t)}{t^{2}}\,\dd t.
\]
代入 \(A(t)=\delta_{\mathrm{ZG}}t+E(t)\) 可得
\[
\sum_{\substack{n\le X\\ n\in\mathrm{ZG}}}\frac{1}{n}
=\delta_{\mathrm{ZG}}\log X+\delta_{\mathrm{ZG}}
+\frac{E(X)}{X}
+\int_{1}^{X}\frac{E(t)}{t^{2}}\,\dd t.
\]
由于
\[
\frac{E(X)}{X}=O_\varepsilon\!\bigl(X^{-\frac12+\varepsilon}\bigr),
\qquad
\frac{E(t)}{t^{2}}=O_\varepsilon\!\bigl(t^{-\frac32+\varepsilon}\bigr),
\]
且指数 \(-\tfrac32+\varepsilon<-1\)，故积分
\[
\int_{1}^{\infty}\frac{E(t)}{t^{2}}\,\dd t
\]
绝对收敛。定义
\[
C_{\mathrm{ZG}}
:=\delta_{\mathrm{ZG}}+\int_{1}^{\infty}\frac{E(t)}{t^{2}}\,\dd t.
\]
于是
\[
\int_{1}^{X}\frac{E(t)}{t^{2}}\,\dd t
=C_{\mathrm{ZG}}-\delta_{\mathrm{ZG}}+O_\varepsilon\!\bigl(X^{-\frac12+\varepsilon}\bigr),
\]
从而得到所述渐近展开。特别地，主项为 \(\delta_{\mathrm{ZG}}\log X\)，故结论成立。
\end{proof}



\begin{theorem}[Zeckendorf--Gödel 素数码集合的自然密度、闭式可计算性与尾项误差界]\label{con:xi-zg-density-closed-form-tail-bound}
对每个 \(N\ge 1\)，定义截断集合
\[
\mathrm{ZG}^{(N)}
:=\Bigl\{
n\in\NN:\ \forall 1\le k\le N,\ v_{p_k}(n)\in\{0,1\},\
\forall 1\le k\le N-1,\ v_{p_k}(n)\,v_{p_{k+1}}(n)=0
\Bigr\}.
\]
记
\[
\delta_N:=\mathrm{dens}\bigl(\mathrm{ZG}^{(N)}\bigr).
\]
则自然密度
\[
\delta_{\mathrm{ZG}}:=\mathrm{dens}(\mathrm{ZG})
\]
存在，且
\[
\delta_{\mathrm{ZG}}=\lim_{N\to\infty}\delta_N\in(0,1).
\]
进一步，\(\delta_N\) 具有显式闭式
\[
\delta_N=\Bigl(\prod_{k=1}^N\Bigl(1-\frac1{p_k}\Bigr)\Bigr)I_N,
\qquad
I_N:=\sum_{\substack{S\subseteq\{1,\dots,N\}\\ S\ \mathrm{无相邻}}}\ \prod_{k\in S}\frac1{p_k},
\]
并满足尾项误差上界
\[
0\le \delta_N-\delta_{\mathrm{ZG}}
\le
\sum_{k>N}\frac1{p_k^2}
\;+\;
\sum_{k\ge N}\frac1{p_kp_{k+1}}.
\]
特别地，\(\delta_{\mathrm{ZG}}\) 可通过生成素数并迭代计算 \(\delta_N\) 在任意精度下有效逼近。
\end{theorem}
\begin{proof}
对固定 \(N\)，成员资格仅依赖于模
\(
M_N:=\prod_{k=1}^N p_k^2
\)
的剩余类，故 \(\delta_N\) 存在。

记
\(
a_k:=1-\frac1{p_k},\ b_k:=a_k\frac1{p_k}
\)。
在前 \(N\) 个素数坐标上，“允许指数为 \(1\)”的索引集合恰为路径图 \(\{1,\dots,N\}\) 的独立集 \(S\)；
对应权重为
\(
\prod_{k\in S}b_k\prod_{k\notin S}a_k
\)。
求和得
\[
\delta_N
=\sum_{S\ \mathrm{无相邻}}\ \prod_{k\in S}b_k\prod_{k\notin S}a_k
=\Bigl(\prod_{k=1}^N a_k\Bigr)\sum_{S\ \mathrm{无相邻}}\prod_{k\in S}\frac1{p_k},
\]
即所示闭式。

又 \(\mathrm{ZG}=\bigcap_{N\ge 1}\mathrm{ZG}^{(N)}\)，故 \(\delta_N\) 单调下降且
\[
\overline{\mathrm{dens}}(\mathrm{ZG})\le \delta_N\quad(\forall N).
\]
若 \(n\in \mathrm{ZG}^{(N)}\setminus \mathrm{ZG}\)，则必发生尾部违约事件之一：
\[
\exists\,k>N,\ p_k^2\mid n
\qquad\text{或}\qquad
\exists\,k\ge N,\ p_kp_{k+1}\mid n.
\]
故
\[
\mathrm{ZG}^{(N)}\setminus \mathrm{ZG}
\subseteq
\Bigl(\bigcup_{k>N}\{p_k^2\mid n\}\Bigr)
\cup
\Bigl(\bigcup_{k\ge N}\{p_kp_{k+1}\mid n\}\Bigr).
\]
由并集上界与
\(
\mathrm{dens}\{m\mid n\}=1/m
\)
得
\[
\mathrm{dens}\bigl(\mathrm{ZG}^{(N)}\setminus \mathrm{ZG}\bigr)
\le
\sum_{k>N}\frac1{p_k^2}
+
\sum_{k\ge N}\frac1{p_kp_{k+1}},
\]
即误差估计。令 \(N\to\infty\) 得上下密度同限，故自然密度存在且等于 \(\lim\delta_N\)。

最后，
\[
\delta_{\mathrm{ZG}}
\ge 1-\sum_{k\ge 1}\frac1{p_k^2}-\sum_{k\ge 1}\frac1{p_kp_{k+1}}>0,
\]
因为两级数收敛且其和严格小于 \(1\)。证毕。
\end{proof}

\begin{corollary}[Mertens 重整常数与素数路径独立集分区函数的对数增长律]\label{cor:xi-zg-mertens-log-growth}
沿用定理 \ref{con:xi-zg-density-closed-form-tail-bound} 的记号。
则
\[
\delta_{\mathrm{ZG}}
=\lim_{N\to\infty}
\Bigl(\prod_{k=1}^N\Bigl(1-\frac1{p_k}\Bigr)\Bigr)I_N,
\]
并且
\[
I_N\sim \delta_{\mathrm{ZG}}\,e^{\gamma}\,\log p_N
\qquad (N\to\infty),
\]
其中 \(\gamma\) 为 Euler--Mascheroni 常数。
\end{corollary}
\begin{proof}
由定理 \ref{con:xi-zg-density-closed-form-tail-bound}，
\[
I_N=\frac{\delta_N}{\prod_{k\le N}(1-\frac1{p_k})}.
\]
再用 Mertens 乘积渐近
\(
\prod_{p\le y}(1-\frac1p)\sim e^{-\gamma}/\log y
\)
并令 \(y=p_N\)、\(\delta_N\to\delta_{\mathrm{ZG}}\) 即得。证毕。
\end{proof}

\begin{theorem}[素数权重 Fibonacci 路径的 Lee--Yang 实根性与 Poisson--binomial 分解]\label{con:xi-zg-weighted-path-leyang-poisson-binomial}
定义
\[
P_N(t):=\sum_{\substack{S\subseteq\{1,\dots,N\}\\ S\ \mathrm{无相邻}}}
t^{|S|}\prod_{k\in S}\frac1{p_k},
\]
则
\[
P_0(t)=1,\qquad
P_1(t)=1+\frac{t}{p_1},\qquad
P_N(t)=P_{N-1}(t)+\frac{t}{p_N}P_{N-2}(t)\quad(N\ge 2).
\]
并且成立：
\begin{enumerate}
  \item \(P_N\) 的全部零点均为互异负实根；
  \item \(P_N\) 与 \(P_{N-1}\) 的零点严格交错；
  \item 设 \(m_N:=\deg P_N=\lfloor(N+1)/2\rfloor\)，则存在 \(\alpha_{N,j}>0\) 使
  \[
  P_N(t)=\prod_{j=1}^{m_N}\Bigl(1+\frac{t}{\alpha_{N,j}}\Bigr).
  \]
\end{enumerate}
对固定 \(t>0\)，令
\[
\pi_{N,j}(t):=\frac{t}{t+\alpha_{N,j}}\in(0,1).
\]
若按 Gibbs 权重
\(
\propto t^{|S|}\prod_{k\in S}p_k^{-1}
\)
在独立集上取样，则随机变量 \(|S|\) 满足分布分解
\[
|S|\overset{d}{=}\sum_{j=1}^{m_N}B_{N,j},
\qquad
B_{N,j}\sim\mathrm{Bernoulli}\!\bigl(\pi_{N,j}(t)\bigr),
\]
其中 \((B_{N,j})\) 相互独立。
\end{theorem}
\begin{proof}
递推直接由索引 \(N\) 是否入独立集的分拆得到。

令 \(\mathcal T_{N+1}\) 为长度 \(N+1\) 的路径树，赋边权
\(
w_k:=p_k^{-1/2}
\)
（\(1\le k\le N\)），并记 \(A_{N+1}\) 为其加权邻接矩阵。
匹配多项式与独立集母函数的标准对应给出
\[
\chi_{N+1}(x):=\det(xI-A_{N+1})
=x^{N+1}P_N\!\Bigl(-\frac1{x^2}\Bigr).
\]
由于 \(A_{N+1}\) 为实对称 Jacobi 矩阵（副对角元严格正），其特征值实且互异；
故 \(P_N\) 的零点为 \(-1/\lambda_j^2\)，均为互异负实数。
交错性由主子矩阵谱交错（Cauchy interlacing）传递到 \(P_N,P_{N-1}\) 的零点。

第三条由实负根分解直接成立。
设
\[
G_{N,t}(s):=\frac{P_N(ts)}{P_N(t)}
=\prod_{j=1}^{m_N}
\frac{1+\frac{ts}{\alpha_{N,j}}}{1+\frac{t}{\alpha_{N,j}}}
=\prod_{j=1}^{m_N}\bigl(1-\pi_{N,j}(t)+\pi_{N,j}(t)s\bigr),
\]
右侧即独立 Bernoulli 和的概率母函数，故得到分布恒等式。证毕。
\end{proof}

\paragraph{素数权硬核归一化常数、Mertens--Lee--Yang 渐近与系数不等式链}
沿用定理 \ref{con:xi-zg-weighted-path-leyang-poisson-binomial} 的素数权路径独立集分区函数 \(P_N(t)\) 与其递推。
\par
\begin{theorem}[硬核归一化常数的存在性与显式尾项控制]\label{thm:xi-prime-hardcore-normalization-constant}
\leanverified{paper\_xi\_prime\_hardcore\_normalization\_constant}
对 \(t>0\) 与 \(N\ge 1\) 定义归一化比值
\begin{equation}\label{eq:xi-prime-hardcore-RN-def}
R_N(t)
:=
\frac{P_N(t)}{\prod_{k=1}^{N}\left(1+\frac{t}{p_k}\right)}\in(0,1].
\end{equation}
则：
\begin{enumerate}
  \item \((R_N(t))_{N\ge 1}\) 严格单调递减；
  \item 极限
  \begin{equation}\label{eq:xi-prime-hardcore-deltaHC-def}
  \boxed{\ \delta_{\mathrm{HC}}(t):=\lim_{N\to\infty}R_N(t)\ }
  \end{equation}
  存在且满足 \(0<\delta_{\mathrm{HC}}(t)<1\)；
  \item 对所有 \(N\ge 2\) 有尾项界
  \begin{equation}\label{eq:xi-prime-hardcore-RN-tail}
  \boxed{\;
  0\ \le\ R_N(t)-\delta_{\mathrm{HC}}(t)\ \le\ \sum_{k\ge N}\frac{t^2}{p_kp_{k+1}}\; }.
  \end{equation}
\end{enumerate}
\end{theorem}
\begin{proof}
记 \(a_k:=t/p_k\) 与 \(D_N(t):=\prod_{k=1}^{N}(1+a_k)\)。
由递推 \(P_N=P_{N-1}+a_NP_{N-2}\) 得
\[
R_N=\frac{R_{N-1}}{1+a_N}+\frac{a_N}{(1+a_{N-1})(1+a_N)}R_{N-2}.
\]
对 \(N\ge 3\) 进一步消元得到差分恒等式
\[
R_{N-1}-R_N
=\frac{a_{N-1}a_N}{(1+a_{N-2})(1+a_{N-1})(1+a_N)}\,R_{N-3}\ >\ 0,
\]
从而 \((R_N)\) 严格递减。
又由于 \(P_N\) 计数的独立集严格包含于所有子集，故 \(0<P_N\le D_N\)，从而 \(0<R_N\le 1\)。
\par
由上式立得上界 \(0<R_{N-1}-R_N\le a_{N-1}a_N=t^2/(p_{N-1}p_N)\)。
由于 \(\sum_{N\ge 2}p_{N-1}^{-1}p_N^{-1}\) 收敛，故 \(\sum_{N\ge 2}(R_{N-1}-R_N)\) 收敛，进而 \(R_N\) 收敛并满足尾项界 \eqref{eq:xi-prime-hardcore-RN-tail}。
\par
最后证明极限严格为正。
由 \(R_{N-2}\ge R_{N-1}\) 与上一差分恒等式得
\[
R_N
\ge R_{N-1}\left(1-\frac{a_{N-1}a_N}{(1+a_{N-1})(1+a_N)}\right),
\]
对 \(N\) 连乘并用 \(\sum_{N\ge 2}a_{N-1}a_N<\infty\) 得无穷乘积下界严格为正，从而 \(\delta_{\mathrm{HC}}(t)>0\)。证毕。
\end{proof}

\begin{corollary}[\texorpdfstring{$t=1$}{t=1} 与 \texorpdfstring{$\mathfrak C_{\mathrm{HC}}$}{C_HC} 的一致性]\label{cor:xi-prime-hardcore-deltaHC-matches-C_HC}
\leanverified{paper\_xi\_prime\_hardcore\_deltahc\_matches\_c\_hc}
在推论 \ref{cor:xi-zg-hardcore-constant-residue} 的记号下，有
\[
\delta_{\mathrm{HC}}(1)=\mathfrak C_{\mathrm{HC}}.
\]
从而
\(
\delta_{\mathrm{ZG}}=\delta_{\mathrm{HC}}(1)/\zeta(2)
\)
与推论 \ref{cor:xi-zg-hardcore-constant-residue} 的闭式一致。
\end{corollary}
\begin{proof}
由 \eqref{eq:xi-prime-hardcore-RN-def} 在 \(t=1\) 的定义，\(\delta_{\mathrm{HC}}(1)=\lim_{N\to\infty}P_N(1)/\prod_{k\le N}(1+1/p_k)\)。
而推论 \ref{cor:xi-zg-hardcore-constant-residue} 中 \(\mathfrak C_{\mathrm{HC}}\) 正是同一极限。证毕。
\end{proof}

\begin{theorem}[Mertens--Lee--Yang 渐近与两常数分解]\label{thm:xi-prime-hardcore-mertens-leyang-asymptotic}
令 \(\delta_{\mathrm{HC}}(t)\) 如定理 \ref{thm:xi-prime-hardcore-normalization-constant} 定义，并定义收敛欧拉乘积
\begin{equation}\label{eq:xi-prime-hardcore-Ct-euler}
C(t):=\prod_{p}\Bigl(1-\frac{1}{p}\Bigr)^{t}\Bigl(1+\frac{t}{p}\Bigr)\in(0,\infty).
\end{equation}
则当 \(N\to\infty\) 时有
\begin{equation}\label{eq:xi-prime-hardcore-asymptotic}
\boxed{\;
P_N(t)\ \sim\ \delta_{\mathrm{HC}}(t)\,C(t)\,e^{\gamma t}\,(\log p_N)^{t}\;}
\end{equation}
其中 \(\gamma\) 为 Euler--Mascheroni 常数。
\end{theorem}
\begin{proof}
由 \eqref{eq:xi-prime-hardcore-RN-def} 与 \eqref{eq:xi-prime-hardcore-deltaHC-def}，
\(
P_N(t)=R_N(t)\prod_{k\le N}(1+t/p_k)\sim \delta_{\mathrm{HC}}(t)\prod_{k\le N}(1+t/p_k)
\)。
将有限乘积分解为
\[
\prod_{k\le N}\Bigl(1+\frac{t}{p_k}\Bigr)
=
\Bigl(\prod_{k\le N}\Bigl(1-\frac{1}{p_k}\Bigr)^{-t}\Bigr)
\cdot
\Bigl(\prod_{k\le N}\Bigl(1-\frac{1}{p_k}\Bigr)^t\Bigl(1+\frac{t}{p_k}\Bigr)\Bigr).
\]
第二个括号内因子满足
\(
\bigl(1-\frac1p\bigr)^t(1+\frac{t}{p})
=1-\frac{t(t+1)}{2p^2}+O(p^{-3})
\)，
其对数为 \(O(p^{-2})\)，故乘积收敛并趋于 \(C(t)\)。
第一个括号由 Mertens 乘积渐近
\(
\prod_{p\le y}(1-\frac1p)\sim e^{-\gamma}/\log y
\)
（取 \(y=p_N\)）给出
\(
\prod_{k\le N}(1-\frac1{p_k})^{-t}\sim e^{\gamma t}(\log p_N)^t
\)。
合并即得 \eqref{eq:xi-prime-hardcore-asymptotic}。证毕。
\end{proof}

\begin{theorem}[素数权硬核 Gibbs 场的规模涨落与中心极限定理]\label{thm:xi-prime-hardcore-gibbs-clt}
\leanverified{paper\_xi\_prime\_hardcore\_gibbs\_clt}
固定 \(t>0\)。
在 \(\{1,\dots,N\}\) 的无相邻集合上按权重
\(
\propto t^{|S|}\prod_{k\in S}p_k^{-1}
\)
抽样，并记 \(K_N:=|S|\)。
则存在仅依赖于 \(t\) 的常数 \(\mu(t),\nu(t)\in\RR\)，使得
\[
\EE[K_N]=t\log\log p_N+\mu(t)+o(1),
\qquad
\mathrm{Var}(K_N)=t\log\log p_N+\nu(t)+o(1).
\]
并且令 \(V_N:=\mathrm{Var}(K_N)\)，则 \(V_N\to\infty\) 且存在绝对常数 \(C>0\) 使
\[
\sup_{x\in\RR}\left|
\PP\!\left(\frac{K_N-\EE[K_N]}{\sqrt{V_N}}\le x\right)-\Phi(x)
\right|
\le \frac{C}{\sqrt{V_N}}
=O\!\left(\frac{1}{\sqrt{\log\log p_N}}\right),
\]
从而 \(\frac{K_N-\EE[K_N]}{\sqrt{V_N}}\Rightarrow\mathcal N(0,1)\)。
\end{theorem}
\begin{proof}
由定理 \ref{thm:xi-prime-hardcore-mertens-leyang-asymptotic}，
\[
\log P_N(t)=t\log\log p_N+\log\delta_{\mathrm{HC}}(t)+\log C(t)+\gamma t+o(1).
\]
对该指数族分配，有标准恒等式
\(
\EE[K_N]=t\,\partial_t\log P_N(t)
\)
与
\(
\mathrm{Var}(K_N)=t\,\partial_t\EE[K_N]
\)，
从而得到所示两条渐近；其中常数 \(\mu(t),\nu(t)\) 由 \(\log\delta_{\mathrm{HC}}(t)+\log C(t)+\gamma t\) 的一、二阶导数给出。
\par
另一方面，由定理 \ref{con:xi-zg-weighted-path-leyang-poisson-binomial} 的 Poisson--binomial 分解，
\(
K_N\overset{d}{=}\sum_{j=1}^{m_N}B_{N,j}
\)
为独立 Bernoulli 和，且 \(\sum_j\mathrm{Var}(B_{N,j})=V_N\to\infty\)。
对 Poisson--binomial 和应用 Berry--Esseen 型界（以第三绝对中心矩上界为代价）得到误差 \(O(V_N^{-1/2})\)，从而得到中心极限定理。证毕。
\end{proof}

\begin{corollary}[单位权素数硬核 Gibbs 场的中心极限定律]\label{cor:xi-prime-hardcore-gibbs-clt-t1-explicit}
\leanverified{paper\_xi\_prime\_hardcore\_gibbs\_clt\_t1\_explicit}
在定理 \ref{thm:xi-prime-hardcore-gibbs-clt} 中取 \(t=1\)。
记
\[
I_N
:=
\sum_{\substack{S\subseteq\{1,\dots,N\}\\ S\ \mathrm{无相邻}}}\prod_{i\in S}\frac1{p_i},
\qquad
\PP_N(S):=\frac{1}{I_N}\prod_{i\in S}\frac1{p_i},
\]
并令 \(\omega_N:=|S|\)。
则
\[
\frac{\omega_N-\mu_N}{\sigma_N}\Longrightarrow \mathcal N(0,1),
\]
其中
\[
\mu_N=\sum_{i=1}^{N}\frac{1}{p_i+1}+O(1),
\qquad
\sigma_N^2=\sum_{i=1}^{N}\frac{p_i}{(p_i+1)^2}+O(1),
\]
并且
\[
\sigma_N^2\sim \sum_{i\le N}\frac1{p_i}\sim \log\log p_N.
\]
\end{corollary}
\begin{proof}
定理 \ref{thm:xi-prime-hardcore-gibbs-clt} 在 \(t=1\) 时给出
\[
\EE[\omega_N]=\log\log p_N+O(1),
\qquad
\mathrm{Var}(\omega_N)=\log\log p_N+O(1),
\]
并且
\(
(\omega_N-\EE[\omega_N])/\sqrt{\mathrm{Var}(\omega_N)}
\Rightarrow
\mathcal N(0,1)
\)。

另一方面，
\[
\frac1{p_i}-\frac{1}{p_i+1}=\frac{1}{p_i(p_i+1)}=O(p_i^{-2}),
\]
以及
\[
\frac1{p_i}-\frac{p_i}{(p_i+1)^2}
=
\frac{2p_i+1}{p_i(p_i+1)^2}
=
O(p_i^{-2}).
\]
由于 \(\sum_i p_i^{-2}<\infty\)，故
\[
\sum_{i=1}^{N}\frac{1}{p_i+1}
=
\sum_{i=1}^{N}\frac1{p_i}+O(1),
\qquad
\sum_{i=1}^{N}\frac{p_i}{(p_i+1)^2}
=
\sum_{i=1}^{N}\frac1{p_i}+O(1).
\]
再由素数倒数和的 Mertens 渐近，
\[
\sum_{i=1}^{N}\frac1{p_i}=\log\log p_N+O(1),
\]
即得所示 \(\mu_N,\sigma_N^2\) 的表达式与发散结论。
最后，由 \(\sigma_N^2=\mathrm{Var}(\omega_N)+O(1)\) 且 \(\sigma_N^2\to\infty\)，两种标准化只相差 \(1+o(1)\) 因子，从而极限正态性保持。证毕。
\end{proof}

\begin{proposition}[素数寄存器与 Lee--Yang 谱的 Newton--Maclaurin 对应及不等式链]\label{prop:xi-prime-hardcore-newton-maclaurin}
\leanverified{paper\_xi\_prime\_hardcore\_newton\_maclaurin}
设 \(m_N=\deg P_N=\lfloor(N+1)/2\rfloor\)，并写
\[
P_N(t)=\sum_{k=0}^{m_N}a_{N,k}\,t^k.
\]
则：
\begin{enumerate}
  \item 系数具有素数寄存器刻画
  \[
  a_{N,k}
  =\sum_{\substack{S\subseteq\{1,\dots,N\}\\ S\ \mathrm{无相邻},\ |S|=k}}
  \ \prod_{i\in S}\frac1{p_i}.
  \]
  \item 若 \(P_N(t)=\prod_{j=1}^{m_N}(1+t/\alpha_{N,j})\)（\(\alpha_{N,j}>0\)；定理 \ref{con:xi-zg-weighted-path-leyang-poisson-binomial}），则
  \[
  a_{N,k}=e_k\!\left(\frac{1}{\alpha_{N,1}},\dots,\frac{1}{\alpha_{N,m_N}}\right),
  \]
  其中 \(e_k\) 为第 \(k\) 个初等对称多项式；并且
  \[
  \sum_{j=1}^{m_N}\frac{1}{\alpha_{N,j}}=\sum_{i=1}^{N}\frac{1}{p_i}.
  \]
  \item 序列
  \[
  \left(\frac{a_{N,k}}{\binom{m_N}{k}}\right)^{1/k}\qquad(k=1,\dots,m_N)
  \]
  随 \(k\) 单调不增；并且 \(a_{N,k}\) 满足严格 Newton 不等式 \(a_{N,k}^2>a_{N,k-1}a_{N,k+1}\)（\(1\le k\le m_N-1\)）。
\end{enumerate}
\end{proposition}
\begin{proof}
前两条由 \(P_N\) 的独立集展开与 Lee--Yang 实根分解展开逐项比较得到。
第三条由正数列的 Maclaurin 不等式与实负互异零点多项式的严格 Newton 不等式推出。证毕。
\end{proof}


\begin{theorem}[Lee--Yang 单位圆提升、Joukowsky 椭圆输运与 \texorpdfstring{$\delta_{\mathrm{ZG}}$}{deltaZG} 的谱行列式表达]\label{con:xi-zg-unit-circle-joukowsky-spectral-determinant}
在定理 \ref{con:xi-zg-weighted-path-leyang-poisson-binomial} 的记号下：
\begin{enumerate}
  \item 设 \(\chi_{N+1}(x)=\det(xI-A_{N+1})\)，\(\Phi_{N+1}(x):=\chi_{N+1}(x/\sqrt2)\)。
  由于 \(\|A_{N+1}\|_{2}\le 2\max_k w_k=\sqrt2\)，\(\Phi_{N+1}\) 全部根落在 \([-2,2]\)。
  因此
  \[
  \mathcal L_{N+1}(z)
  :=z^{N+1}\Phi_{N+1}(z+z^{-1})
  =z^{N+1}\chi_{N+1}\!\Bigl(\frac{z+z^{-1}}{\sqrt2}\Bigr)
  \]
  的全部零点位于单位圆 \(|z|=1\)。
  \item 对任意 \(r>1\)，令 \(J_r(z):=rz+r^{-1}z^{-1}\)。
  若 \(z_j\) 为 \(\mathcal L_{N+1}\) 的零点，则 \(w_j:=J_r(z_j)\) 落在椭圆 \(J_r(S^1)\) 上。
  \item 取 \(x=i\) 于恒等式
  \(
  \chi_{N+1}(x)=x^{N+1}P_N(-x^{-2})
  \)
  得
  \[
  P_N(1)=|\chi_{N+1}(i)|=|\det(iI-A_{N+1})|.
  \]
  因而
  \[
  \delta_{\mathrm{ZG}}
  =\lim_{N\to\infty}
  \Bigl(\prod_{k=1}^N\Bigl(1-\frac1{p_k}\Bigr)\Bigr)\,|\det(iI-A_{N+1})|.
  \]
  又
  \[
  |\det(iI-A_{N+1})|^2
  =\det\!\bigl((iI-A_{N+1})(-iI-A_{N+1})\bigr)
  =\det(I+A_{N+1}^2),
  \]
  故
  \[
  P_N(1)=\det(I+A_{N+1}^2)^{1/2},
  \quad
  \delta_{\mathrm{ZG}}
  =\lim_{N\to\infty}
  \Bigl(\prod_{k=1}^N\Bigl(1-\frac1{p_k}\Bigr)\Bigr)\det(I+A_{N+1}^2)^{1/2}.
  \]
\end{enumerate}
\end{theorem}
\begin{proof}
第一条由谱半径界和参数化 \(x=z+z^{-1}=2\cos\theta\) 直接得到。
第二条是 \(J_r\) 对单位圆的标准像性质。
第三条由 \(x=i\) 代入与行列式恒等式逐步化简得出。证毕。
\end{proof}

\begin{corollary}[\texorpdfstring{$s=1$}{s=1} 处单极点与 Dirichlet 留数常数]\label{cor:xi-zg-dirichlet-residue}
令 \(P(s)=\sum_{k\ge 1}p_k^{-s}\) 为 prime zeta，并令 \(G(s)\) 为定理 \ref{thm:xi-zg-dirichlet-log-renormalization} 的正则因子。
采用 prime zeta 在 \(s\to 1^+\) 的经典展开
\[
P(s)=\log\frac1{s-1}+B_{\mathsf P}+o(1),
\]
其中 \(B_{\mathsf P}\in\RR\) 为常数（可由 Mertens 定理推导；参见 \cite{Apostol1976,Titchmarsh1986RiemannZeta}）。
则 \(F(s)=\sum_{n\in\mathrm{ZG}}n^{-s}\) 在 \(s=1\) 处具有单极点，且其留数为
\[
\lim_{s\to1^+}(s-1)F(s)=\exp\!\bigl(B_{\mathsf P}+G(1)\bigr)\in(0,\infty).
\]
并且该留数等于 \(\mathrm{ZG}\) 的 Dirichlet 密度；结合定理 \ref{con:xi-zg-density-closed-form-tail-bound} 的自然密度存在性，两者一致。
\end{corollary}

\paragraph{Abel 边界律与前缀递推证书}

\begin{theorem}[ZG 的 Dirichlet 级数在 \texorpdfstring{$s=1$}{s=1} 的 Abel 型边界律、调和主项与收敛临界线]\label{thm:xi-zg-abel-residue-log-density}
沿用本段记号，记
\[
A_{\mathrm{ZG}}(x):=\#\{n\le x:\ n\in\mathrm{ZG}\},
\qquad
D_{\mathrm{ZG}}(s):=\sum_{n\in\mathrm{ZG}}n^{-s}
\qquad (\Re(s)>1),
\]
以及
\[
\delta_{\mathrm{ZG}}:=\lim_{x\to\infty}\frac{A_{\mathrm{ZG}}(x)}{x}.
\]
则 \(\delta_{\mathrm{ZG}}\in(0,1)\)，并且成立：
\begin{enumerate}
  \item Abel 型边界律
  \[
  \lim_{\sigma\downarrow 1}(\sigma-1)D_{\mathrm{ZG}}(\sigma)=\delta_{\mathrm{ZG}}.
  \]
  \item 调和和主项
  \[
  \sum_{\substack{n\le X\\ n\in\mathrm{ZG}}}\frac{1}{n}
  =\delta_{\mathrm{ZG}}\log X+o(\log X)
  \qquad (X\to\infty).
  \]
  \item 绝对收敛临界线满足
  \[
  \sigma_c(D_{\mathrm{ZG}})=1.
  \]
\end{enumerate}
\end{theorem}
\begin{proof}
由定理 \ref{con:xi-zg-density-closed-form-tail-bound}，自然密度存在且
\[
\delta_{\mathrm{ZG}}\in(0,1).
\]
记
\[
A_{\mathrm{ZG}}(x)=\delta_{\mathrm{ZG}}\,x+r(x),
\qquad
r(x)=o(x)
\qquad (x\to\infty).
\]

对任意 \(\sigma>1\)，Abel--Stieltjes 分部积分给出
\[
D_{\mathrm{ZG}}(\sigma)
=
\sigma\int_1^\infty A_{\mathrm{ZG}}(x)x^{-\sigma-1}\,\dd x.
\]
代入 \(A_{\mathrm{ZG}}(x)=\delta_{\mathrm{ZG}}x+r(x)\) 后得到
\[
D_{\mathrm{ZG}}(\sigma)
=
\frac{\sigma\,\delta_{\mathrm{ZG}}}{\sigma-1}
+
\sigma\int_1^\infty r(x)x^{-\sigma-1}\,\dd x.
\]
因此只需证明
\[
\lim_{\sigma\downarrow1}
(\sigma-1)\int_1^\infty r(x)x^{-\sigma-1}\,\dd x=0.
\]
任取 \(\varepsilon>0\)。
由 \(r(x)=o(x)\)，可取 \(X_0\ge 1\) 使得 \(|r(x)|\le \varepsilon x\) 对所有 \(x\ge X_0\) 成立。
于是
\[
\begin{aligned}
(\sigma-1)\left|\sigma\int_1^\infty r(x)x^{-\sigma-1}\,\dd x\right|
&\le
(\sigma-1)\sigma\int_1^{X_0}|r(x)|x^{-\sigma-1}\,\dd x \\
&\qquad
+(\sigma-1)\sigma\int_{X_0}^{\infty}\varepsilon x\cdot x^{-\sigma-1}\,\dd x \\
&\le
(\sigma-1)C_{X_0}
+\sigma\varepsilon X_0^{1-\sigma},
\end{aligned}
\]
其中 \(C_{X_0}<\infty\) 与 \(\sigma\) 无关。
令 \(\sigma\downarrow1\)，上式上极限不超过 \(\varepsilon\)；再令 \(\varepsilon\downarrow0\)，即得
\[
\lim_{\sigma\downarrow 1}(\sigma-1)D_{\mathrm{ZG}}(\sigma)=\delta_{\mathrm{ZG}}.
\]

再看第二条。
对和式作 Abel 分部积分，
\[
\sum_{\substack{n\le X\\ n\in\mathrm{ZG}}}\frac1n
=
\frac{A_{\mathrm{ZG}}(X)}{X}+\int_1^X \frac{A_{\mathrm{ZG}}(t)}{t^2}\,\dd t.
\]
代入 \(A_{\mathrm{ZG}}(t)=\delta_{\mathrm{ZG}}t+r(t)\) 得
\[
\sum_{\substack{n\le X\\ n\in\mathrm{ZG}}}\frac1n
=
\delta_{\mathrm{ZG}}
+\frac{r(X)}{X}
+\delta_{\mathrm{ZG}}\log X
+\int_1^X \frac{r(t)}{t^2}\,\dd t.
\]
由 \(r(t)=o(t)\)，对任意 \(\varepsilon>0\) 仍可取 \(T_0\) 使 \(|r(t)|\le \varepsilon t\) 对 \(t\ge T_0\) 成立，从而
\[
\left|\int_1^X \frac{r(t)}{t^2}\,\dd t\right|
\le C_{T_0}+\varepsilon\log X,
\]
其中 \(C_{T_0}\) 与 \(X\) 无关。
因此
\[
\sum_{\substack{n\le X\\ n\in\mathrm{ZG}}}\frac1n
=
\delta_{\mathrm{ZG}}\log X+o(\log X).
\]
最后，由 \(\mathrm{ZG}\subset \NN\) 可知对任意 \(\sigma>1\) 都有
\[
0\le D_{\mathrm{ZG}}(\sigma)\le \zeta(\sigma)<\infty,
\]
故 \(D_{\mathrm{ZG}}\) 在半平面 \(\Re(s)>1\) 上绝对收敛，从而 \(\sigma_c(D_{\mathrm{ZG}})\le 1\)。另一方面，第二条等价于
\[
\sum_{\substack{n\le X\\ n\in\mathrm{ZG}}}\frac1n
=
\delta_{\mathrm{ZG}}\log X+o(\log X)\to\infty,
\]
从而级数 \(D_{\mathrm{ZG}}(1)=\sum_{n\in\mathrm{ZG}}n^{-1}\) 发散。于是其收敛 abscissa 恰为 \(1\)。
\end{proof}

\begin{corollary}[ZG 的 Dirichlet 级数的绝对收敛临界线恰为 \texorpdfstring{$\Re(s)=1$}{Re(s)=1}]\label{cor:xi-zg-absolute-convergence-critical-line}
在定理 \ref{thm:xi-zg-abel-residue-log-density} 的设定下，\(D_{\mathrm{ZG}}(s)\) 在 \(\Re(s)>1\) 上绝对收敛，在 \(s=1\) 发散，并满足
\[
D_{\mathrm{ZG}}(\sigma)\sim \frac{\delta_{\mathrm{ZG}}}{\sigma-1}
\qquad (\sigma\downarrow 1).
\]
因而 \(\Re(s)=1\) 正是其绝对收敛临界线。
\end{corollary}
\begin{proof}
由定理 \ref{thm:xi-zg-abel-residue-log-density} 的第一条，
\[
(\sigma-1)D_{\mathrm{ZG}}(\sigma)\longrightarrow \delta_{\mathrm{ZG}}>0
\qquad (\sigma\downarrow 1),
\]
故
\[
D_{\mathrm{ZG}}(\sigma)
=
\frac{\delta_{\mathrm{ZG}}}{\sigma-1}
+o\!\bigl((\sigma-1)^{-1}\bigr)
\qquad (\sigma\downarrow 1).
\]
特别地，\(D_{\mathrm{ZG}}(\sigma)\to+\infty\)。再结合同一定理的第三条，便得到所述绝对收敛临界线结论。证毕。
\end{proof}

\begin{proposition}[Zeckendorf--Gödel 密度的归一化前缀递推、单调收敛与尾误差证书]\label{prop:xi-zg-prefix-recursion-monotone-tail-certificate}
沿用定理 \ref{con:xi-zg-density-closed-form-tail-bound} 与
\ref{con:xi-zg-weighted-path-leyang-poisson-binomial} 的记号。
定义数列
\[
(a_0,b_0)=(1,0),
\]
并对 \(N\ge 1\) 递推地令
\[
a_N=(a_{N-1}+b_{N-1})\frac{p_N}{p_N+1},
\qquad
b_N=\frac{a_{N-1}}{p_N+1}.
\]
再记
\[
u_N:=a_N+b_N,
\qquad
\widetilde\delta_N:=\frac{u_N}{\zeta(2)}.
\]
则成立：
\begin{enumerate}
  \item 对每个 \(N\ge 0\)，有
  \[
  u_N=
  \frac{P_N(1)}{\prod_{k=1}^{N}\left(1+\frac1{p_k}\right)},
  \]
  其中约定空积等于 \(1\)。
  因而
  \[
  \delta_{\mathrm{ZG}}
  =
  \frac{1}{\zeta(2)}\lim_{N\to\infty}u_N.
  \]
  \item \(u_0=u_1=1\)，并且数列 \((\widetilde\delta_N)_{N\ge 1}\) 严格单调下降，满足
  \[
  \widetilde\delta_N\downarrow \delta_{\mathrm{ZG}}.
  \]
  \item 对每个 \(N\ge 1\)，有显式尾误差界
  \[
  0\le \widetilde\delta_N-\delta_{\mathrm{ZG}}
  \le
  \widetilde\delta_N
  \sum_{j=N}^{\infty}\frac{1}{(p_j+1)(p_{j+1}+1)}.
  \]
\end{enumerate}
\end{proposition}
\begin{proof}
对 \(N\ge 1\)，定义两类部分配分函数
\[
A_N:=
\sum_{\substack{S\subseteq\{1,\dots,N\}\\ S\ \mathrm{无相邻}\\ N\notin S}}
\prod_{k\in S}\frac1{p_k},
\qquad
B_N:=
\sum_{\substack{S\subseteq\{1,\dots,N\}\\ S\ \mathrm{无相邻}\\ N\in S}}
\prod_{k\in S}\frac1{p_k}.
\]
则
\[
P_N(1)=A_N+B_N.
\]
并且由末位是否被选中，直接得到
\[
A_N=P_{N-1}(1),
\qquad
B_N=\frac{1}{p_N}A_{N-1}.
\]
把它们按
\[
a_N:=\frac{A_N}{\prod_{k=1}^{N}\left(1+\frac1{p_k}\right)},
\qquad
b_N:=\frac{B_N}{\prod_{k=1}^{N}\left(1+\frac1{p_k}\right)}
\]
归一化，便得到所述递推关系
\[
a_N=(a_{N-1}+b_{N-1})\frac{p_N}{p_N+1},
\qquad
b_N=\frac{a_{N-1}}{p_N+1}.
\]
于是
\[
u_N=a_N+b_N
=
\frac{P_N(1)}{\prod_{k=1}^{N}\left(1+\frac1{p_k}\right)}.
\]
这就证明了第一条中的第一式。

另一方面，由定理 \ref{con:xi-zg-density-closed-form-tail-bound}，
\[
\delta_N
=
\left(\prod_{k=1}^{N}\left(1-\frac1{p_k}\right)\right)P_N(1).
\]
将上一式代入并整理得
\[
\delta_N
=
\left(\prod_{k=1}^{N}\left(1-\frac1{p_k^2}\right)\right)u_N.
\]
由于 \(\delta_N\to \delta_{\mathrm{ZG}}\) 且
\[
\prod_{k=1}^{N}\left(1-\frac1{p_k^2}\right)\longrightarrow \frac1{\zeta(2)},
\]
只要 \(u_N\) 收敛，便有
\[
\delta_{\mathrm{ZG}}=\frac{1}{\zeta(2)}\lim_{N\to\infty}u_N.
\]

现证单调性。
由递推直接计算
\[
u_N
=
u_{N-1}-\frac{b_{N-1}}{p_N+1}.
\]
因此 \(u_0=u_1=1\)。
而对 \(N\ge 2\)，由 \(a_{N-2}>0\) 可得
\[
b_{N-1}=\frac{a_{N-2}}{p_{N-1}+1}>0,
\]
故
\[
u_N<u_{N-1}\qquad (N\ge 2).
\]
于是 \((u_N)_{N\ge 1}\) 严格递减且下有界于 \(0\)，故极限
\(
u_\infty:=\lim_{N\to\infty}u_N
\)
存在。
由前述关系，便得到
\[
\widetilde\delta_N=\frac{u_N}{\zeta(2)}\downarrow \frac{u_\infty}{\zeta(2)}=\delta_{\mathrm{ZG}}.
\]

最后证明尾误差界。
由
\[
u_N-u_{N+1}=\frac{b_N}{p_{N+1}+1}
\]
以及
\[
b_N=\frac{a_{N-1}}{p_N+1},
\qquad
u_N=a_{N-1}+\frac{p_N}{p_N+1}b_{N-1}\ge a_{N-1},
\]
可得
\[
u_N-u_{N+1}
\le
\frac{u_N}{(p_N+1)(p_{N+1}+1)}.
\]
而由于 \((u_N)\) 递减，对任意 \(j\ge N\) 亦有
\[
u_j-u_{j+1}
\le
\frac{u_j}{(p_j+1)(p_{j+1}+1)}
\le
\frac{u_N}{(p_j+1)(p_{j+1}+1)}.
\]
故
\[
u_N-u_\infty
=
\sum_{j=N}^{\infty}(u_j-u_{j+1})
\le
u_N\sum_{j=N}^{\infty}\frac{1}{(p_j+1)(p_{j+1}+1)}.
\]
两边乘以 \(\zeta(2)^{-1}\) 即得
\[
0\le \widetilde\delta_N-\delta_{\mathrm{ZG}}
\le
\widetilde\delta_N\sum_{j=N}^{\infty}\frac{1}{(p_j+1)(p_{j+1}+1)}.
\]
证毕。
\end{proof}


\begin{proposition}[临界重整化常数的可计算误差界]\label{prop:xi-zg-holographic-constant-tail-bound}
\leanverified{paper\_xi\_zg\_holographic\_constant\_tail\_bound}
取 \(s=1\)，并令 \(\Delta_N(1)\) 与 \(G(1)=\sum_{N\ge 2}\Delta_N(1)\) 如定理 \ref{thm:xi-zg-dirichlet-log-renormalization}。
定义截断
\[
G^{(N)}(1):=\sum_{k=2}^{N}\Delta_k(1).
\]
则存在绝对常数 \(C>0\) 使对所有 \(N\ge 3\) 有
\[
\bigl|G(1)-G^{(N)}(1)\bigr|
\le
C\sum_{k>N}\Bigl(\frac1{p_k^2}+\frac1{p_kp_{k-1}}\Bigr),
\]
从而正则因子 \(\exp(G(1))\) 可被一条显式可求和尾项在任意精度下有效逼近。
\end{proposition}

\begin{proof}
由定理 \ref{thm:xi-zg-dirichlet-log-renormalization} 的一致估计，在 \(s=1\) 时存在常数 \(C\) 使
\(
|\Delta_k(1)|
\le C\bigl(p_k^{-2}+(p_kp_{k-1})^{-1}\bigr)
\)。
对尾和求并界即得。证毕。
\end{proof}

\begin{theorem}[\(\mathrm{ZG}\) 的对数 Mertens 定律]\label{thm:xi-zg-log-mertens}
设
\[
H(x):=\sum_{\substack{n\le x\\ n\in\mathrm{ZG}}}\frac1n.
\]
则极限存在并满足
\[
\lim_{x\to\infty}\frac{H(x)}{\log x}=\delta_{\mathrm{ZG}},
\]
并且常数项可由 \(F(s)=\sum_{n\in\mathrm{ZG}}n^{-s}\) 在 \(s=1\) 的 Laurent 常数项显式识别：存在 \(C_{\mathrm{ZG}}\in\RR\) 使
\[
H(x)=\delta_{\mathrm{ZG}}\log x+C_{\mathrm{ZG}}+o(1)\qquad(x\to\infty).
\]
其中
\[
C_{\mathrm{ZG}}
=\frac{\mathfrak C_{\mathrm{HC}}}{\zeta(2)}
\left(\gamma-\frac{2\zeta'(2)}{\zeta(2)}+\kappa_{\mathrm{HC}}\right)
=\delta_{\mathrm{ZG}}
\left(\gamma-\frac{2\zeta'(2)}{\zeta(2)}+\kappa_{\mathrm{HC}}\right).
\]
\end{theorem}

\begin{proof}
由推论 \ref{cor:xi-zg-dirichlet-residue} 或推论 \ref{cor:xi-zg-hardcore-constant-residue}，\(F(s)=\sum_{n\in\mathrm{ZG}}n^{-s}\) 在 \(s=1\) 处的留数为 \(\delta_{\mathrm{ZG}}\)。
对非负系数 Dirichlet 级数应用标准 Tauberian 定理可得
\(
H(x)=\delta_{\mathrm{ZG}}\log x+C_{\mathrm{ZG}}+o(1)
\)，其中 \(C_{\mathrm{ZG}}\) 等于 \(F\) 在 \(s=1\) 的 Laurent 常数项。
而命题 \ref{prop:xi-zg-hardcore-laurent} 给出了该 Laurent 常数项的显式表达式，从而得到所示恒等式。证毕。
\end{proof}

\begin{conjecture}[正则因子的全半平面延拓与 Lee--Yang 型零点边界]\label{conj:xi-zg-leyang-godel-transport-boundary}
令 \(\mathcal R(s):=\exp(G(s))\)（定理 \ref{thm:xi-zg-dirichlet-log-renormalization}），其在 \(\Re(s)>\tfrac12\) 全纯且无零。
猜想 \(\mathcal R\) 可延拓为 \(\Re(s)>0\) 上的全纯无零函数；从而 \(F(s)=\exp(P(s))\mathcal R(s)\) 在 \(\Re(s)>0\) 的全部奇性结构由 prime zeta \(P(s)\) 的奇性传输唯一决定。
进一步地，对充分大的截断 \(N\)，将 \(F_N(s)\) 视为有限体积分配函数，则其归一化零点云在 \(N\to\infty\) 的极限中形成一条由 Joukowsky 椭圆输运诱导的边界曲线，并以 \(\Re(s)=\tfrac12\) 为临界对称轴，与 \(\mathcal R\) 的零点真空相容。
\end{conjecture}

\begin{conjecture}[弱局域约束下的 Erd\H{o}s--Kac 型高斯极限]\label{conj:xi-zg-erdos-kac-under-local-hardcore}
记
\[
\omega(n):=\#\{p:\ p\mid n\},
\qquad
\mathcal Z(X):=\{n\le X:\ n\in\mathrm{ZG}\}.
\]
猜想存在常数 \(c_0\in\RR\)，使得在 \(\mathcal Z(X)\) 上均匀抽样的 \(n\) 满足
\[
\frac{\omega(n)-\log\log X-c_0}{\sqrt{\log\log X}}
\Longrightarrow
\mathcal N(0,1)
\qquad(X\to\infty).
\]
该猜想表明：相邻素数不可同现的局域硬核约束在中心极限定律尺度上仅导致有限平移，而不改变标准化极限分布。
\end{conjecture}

\begin{remark}[可复现核验脚本]\label{rem:xi-zg-density-leyang-spectral-audit-script}
定理 \ref{con:xi-zg-density-closed-form-tail-bound}--\ref{con:xi-zg-unit-circle-joukowsky-spectral-determinant}
中的闭式恒等式与有限窗数值核验由脚本
\texttt{scripts/exp\_xi\_prime\_register\_zg\_density\_leyang\_spectrum\_audit.py}
给出，对应证书文件为
\texttt{artifacts/export/xi\_prime\_register\_zg\_density\_leyang\_spectrum\_audit.json}。
\par
定理 \ref{thm:xi-zg-dirichlet-log-renormalization} 与命题 \ref{prop:xi-zg-holographic-constant-tail-bound}
在 \(s=1\) 处的递推与尾项模板核验由脚本
\texttt{scripts/exp\_xi\_prime\_register\_zg\_dirichlet\_log\_renormalization\_audit.py}
给出，对应证书文件为
\texttt{artifacts/export/xi\_prime\_register\_zg\_dirichlet\_log\_renormalization\_audit.json}。
\par
定理 \ref{thm:xi-zg-zeta-factorization}--命题 \ref{prop:xi-zg-hardcore-laurent}
中 \(\mathfrak C_{\mathrm{HC}}\) 与 \(\kappa_{\mathrm{HC}}\) 的数值核验由脚本
\texttt{scripts/exp\_xi\_prime\_register\_zg\_hardcore\_zeta\_factorization\_constant\_audit.py}
给出，对应证书文件为
\texttt{artifacts/export/xi\_prime\_register\_zg\_hardcore\_zeta\_factorization\_constant\_audit.json}。
\end{remark}
